\begin{abstract}
\noindent\textbf{Background.} Large language models (LLMs) can populate structured clinical forms from free
text, but long forms force a design choice that is rarely evaluated directly: request every variable in one
generation, or decompose the form into several calls whose outputs are validated and merged. Independent calls
shorten outputs but lose a common reading of the clinical episode, which many variables require.

\noindent\textbf{Objective.} To compare, on the same patients, models and 199-variable intensive care unit
(ICU) form, three strategies: (A) a single monolithic call; (B) a schema-constrained harness of 18 independent
calls of 2 to 16 variables each; and (C) the same harness in which a preliminary call builds an
evidence-grounded model of the case (ICU stays, timeline, problem status, support, infections, outcome) shared
by all group calls. Outcomes are structural validity, evidence fidelity, cross-block clinical coherence,
agreement with human annotators and cost.

\noindent\textbf{Methods.} Retrospective methodological study on a frozen cohort of 50 de-identified Spanish
discharge summaries from the adult ICU of a public hospital in Chile. All strategies receive the full
line-indexed note, the same annotation rules and the same output contract; evidence is returned as line
identifiers and the quoted text is copied by software. Each strategy is run with three open-weight models on
local infrastructure (Gemma~4 31B, Llama~3.3 70B and Qwen~3.6 35B-A3B) with greedy decoding. The primary
comparison is non-inferiority of B versus A in agreement with a pre-defined human reference in summaries with a
documented ICU stay; C versus B and C versus A are secondary. Differences are estimated within model with a
patient-level paired bootstrap.

\noindent\textbf{Results.} The cohort comprises 50 summaries (median 3 pages and approximately 6{,}000
characters); 12 did not mention an ICU stay and are candidates for exclusion from agreement analyses. At the
23 September 2026 data cut, 63 of 150 planned human annotations had been submitted by 4 reviewers; pairwise
agreement on the 174 boolean variables was $\kappa=0.76$ (95\% CI 0.70--0.82; 19 summaries). In a single-case
pilot with Gemma~4, both A and B returned all 199 variables without truncation; B took 39~min and left 6
relation errors between independently answered groups, and A took 3~h~47~min and was valid after one retry.
Comparative results are \pending{pending the evaluation run}.

\noindent\textbf{Conclusions.} \pending{to be written after the pre-specified analysis. The conclusion must be
limited to the measured paired differences, this cohort, this schema and these three models, and must not
imply clinical use.}

\medskip\noindent\textbf{Keywords:} clinical information extraction; large language models; prompt
engineering; task decomposition; intermediate representation; structured output; discharge summaries;
intensive care; Spanish; evidence grounding; reproducibility
\end{abstract}

\begin{tcolorbox}[colback=gray!5,colframe=gray!50,title=Alternative titles (for author selection),fonttitle=\sffamily\bfseries]
\begin{enumerate}
\item One call, eighteen, or eighteen with a shared case model? A paired evaluation of LLM extraction of a
199-variable ICU form from Spanish discharge summaries.
\item Monolithic, decomposed and case-model-guided extraction of ICU variables with open-weight LLMs: structure,
evidence, coherence and cost.
\item Does sharing a model of the episode make decomposed LLM extraction coherent? A paired study on Spanish ICU
discharge summaries.
\item Auditable extraction of 199 ICU variables with local LLMs: a within-patient comparison of three prompting
strategies.
\item From one long answer to many short ones: decomposition, shared context and cost in LLM form filling from
ICU discharge summaries.
\end{enumerate}
\end{tcolorbox}
